% Preámbulo común de CP4_Bales.pdf. Lo carga main.tex; los fragmentos de cada
% inciso no deben tener preámbulo propio ni cargar paquetes.

\usepackage[T1]{fontenc}
\usepackage[utf8]{inputenc}
\usepackage[spanish,es-tabla,es-nodecimaldot]{babel}
\usepackage{lmodern}
\usepackage{microtype}
\usepackage[a4paper,margin=2.5cm]{geometry}

\usepackage{graphicx}
% Las rutas de las imágenes son relativas a la raíz del repo, por ejemplo:
% \includegraphics{ejercicio-12-transacciones-bancarias/diagramas/mi-diagrama.png}
\graphicspath{{../}}

\usepackage{booktabs}
\usepackage{tabularx}
\usepackage{float}
\usepackage{caption}
\usepackage{enumitem}
\usepackage{amsmath}
\usepackage[table]{xcolor}
\usepackage{tikz}

% ---------- Colores de la portada (misma paleta que los diagramas) ----------
\definecolor{portadaoscuro}{HTML}{1E3A6E}
\definecolor{portadaacento}{HTML}{3D5FB0}
\definecolor{portadaclaro}{HTML}{C9D8F5}
\definecolor{portadagris}{HTML}{6B7280}
\definecolor{portadafila}{HTML}{EEF3FC}
\usepackage{listings}
\usepackage[hidelinks]{hyperref}

% ---------- Código C# ----------
\definecolor{codigofondo}{HTML}{F6F8FA}
\definecolor{codigoclave}{HTML}{0033B3}
\definecolor{codigocomentario}{HTML}{6A737D}
\definecolor{codigocadena}{HTML}{067D17}

\lstdefinestyle{csharp}{
  language=[Sharp]C,
  basicstyle=\ttfamily\small,
  keywordstyle=\color{codigoclave}\bfseries,
  commentstyle=\color{codigocomentario}\itshape,
  stringstyle=\color{codigocadena},
  backgroundcolor=\color{codigofondo},
  frame=single, rulecolor=\color{black!15},
  numbers=left, numberstyle=\tiny\color{black!50},
  breaklines=true, showstringspaces=false, tabsize=4,
  morekeywords={var, async, await, record, init, get, set, required},
  % listings no entiende UTF-8: estas sustituciones permiten tildes y ñ en el código
  literate={á}{{\'a}}1 {é}{{\'e}}1 {í}{{\'i}}1 {ó}{{\'o}}1 {ú}{{\'u}}1
           {Á}{{\'A}}1 {É}{{\'E}}1 {Í}{{\'I}}1 {Ó}{{\'O}}1 {Ú}{{\'U}}1
           {ñ}{{\~n}}1 {Ñ}{{\~N}}1 {ü}{{\"u}}1 {¿}{{?`}}1 {¡}{{!`}}1
           {→}{{$\rightarrow$}}2 {×}{{$\times$}}1
}
\lstset{style=csharp}

% ---------- Incisos ----------
% \inciso{título}{ruta del fragmento desde la raíz del repo}
% main.tex pone el título de la subsección; el fragmento solo trae el contenido.
% Si el fragmento todavía no existe, se muestra "Pendiente" y el PDF sigue compilando.
\newcommand{\inciso}[2]{%
  \subsection{#1}%
  \InputIfFileExists{../#2}{}{\textit{Pendiente.}}%
}
